\documentclass[11pt,a4paper]{article}
\usepackage[T1]{fontenc}
\usepackage[utf8]{inputenc}
\usepackage{amsmath,amssymb,amsthm}
\usepackage[margin=1in]{geometry}
\usepackage[colorlinks=true,linkcolor=blue,urlcolor=blue]{hyperref}
\theoremstyle{plain}
\newtheorem{theorem}{Theorem}
\newtheorem{lemma}[theorem]{Lemma}
\newtheorem{proposition}[theorem]{Proposition}
\newtheorem{corollary}[theorem]{Corollary}
\theoremstyle{definition}
\newtheorem{remark}[theorem]{Remark}

\title{The empirical recurrence for OEIS A183928, proved from an exact model}
\author{Adrian Perez Fontelles\\ \small Independent researcher}
\date{13 September 2026}
\begin{document}
\maketitle

\begin{abstract}
OEIS A183928 carries an empirical recurrence of order $11$. It is true, and it is
decidable rather than empirical. The entry's sequence is given exactly by an Ehrhart quasi-polynomial of the multiset cone, from
which $a$ provably satisfies a MONIC linear recurrence of order $S=13$, derived below and not
assumed. The residual of the conjectured recurrence satisfies the same one, so whether it
annihilates $a$ is settled by a finite exact computation. No transfer matrix is involved and
the count is not a walk.
\end{abstract}

\noindent\small 2020 Mathematics Subject Classification. 05A15, 05A19, 11P21.\normalsize

\section{The conjecture}

OEIS A183928 is ``Number of nondecreasing arrangements of n numbers in -2..2 with sum zero and sum of squares less than 2n.''. It has offset $1$ and begins
\[
1, 2, 2, 5, 8, 10, 15, 21,\ \dots
\]
The entry states, as a conjecture contributed by its author and never marked settled:
\begin{quote}
Empirical: a(n) = a(n-1)+2*a(n-3)-a(n-4)-a(n-5)-a(n-6)-a(n-7)+2*a(n-8)+a(n-10)-a(n-11).
\end{quote}
As of the ``Last modified'' line on the live entry (Apr 11 2018 23:41:31, revision 36)
this is still recorded as empirical, and nothing on the entry records it as proved.

\section{The count is an Ehrhart quasi-polynomial}

A nondecreasing arrangement of $n$ numbers from $-2..2$ is a MULTISET, so it is its counts
$c_v$ for $v=-2,\dots,2$, and every condition the entry names is linear in those and in $n$:
\[
\sum_v c_v = n,\qquad \sum_v v\,c_v = 0,\qquad c_v\ge0 and $\sum_v v^2 c_v < 2 n$.
\]
Nothing here is a walk: the LENGTH grows and the alphabet is fixed, so no window counts it. Use
the first equation to eliminate $n$; every remaining condition is then homogeneous in $c$ alone,
the admissible $c$ form a finite union of relatively open rational cones in $\mathbf{R}^{2\cdot2+1}$,
and $a(n)$ counts their lattice points at height $\sum_v c_v = n$. It is therefore an Ehrhart
quasi-polynomial, and the count itself is an exact integer dynamic programme over $v$ carrying
the running sums.

\section{The bound}

The period is read off the arrangement rather than assumed. Every ray of every cell is cut out
by as many independent hyperplanes as there are variables less one, so by Cramer its primitive
generator is the vector of cofactors, and its height $\sum_v c_v$ is what the period divides. A
ray leaving the region --- some $c_v$ negative, or the wrong side of one of the two other
conditions --- bounds no cell and is dropped. Writing $T$ for the heights that survive, the
annihilator is $\prod_{d\,\mid\,t\in T}\Phi_d(z)^{m_d}$ with $m_d$ at most the dimension of the
cone, which the equality $\sum_v v c_v=0$ drops to $4$; one further factor $z$ covers the
origin, the one cell with no ray, whose series is the constant $1$. That gives a monic
annihilator of order $S=13$, which was then asked for terms the model had not supplied
and reproduced them.

\section{The proof}

\begin{lemma}\label{lem:crit}
Let $A(z)=z^S-\sum_{j<S}\alpha_jz^{j}$ be the monic annihilator derived above and
$q(z)=z^{11}-\sum_ic_iz^{11-i}$ the characteristic polynomial of the conjectured
recurrence. The residual $r(n)=a(n)-\sum_ic_ia(n-i)$ is a linear combination of shifts of $a$,
so $A$ annihilates $r$ as well. Since $A(0)\ne0$ the recurrence it gives runs in both
directions, and $S$ consecutive zeros of $r$ force $r$ to vanish at every later index.
\end{lemma}

\begin{proof}
$A$ annihilates $a$, hence every shift of $a$, hence any linear combination of shifts; that is
$r$. A monic recurrence with nonzero constant term determines each term from the $S$ before it
and each from the $S$ after it, so a run of $S$ zeros propagates.
\end{proof}

\begin{theorem}
\[
a(n)\;=\;a(n-1) + 2\,a(n-3) - a(n-4) - a(n-5) - a(n-6) - a(n-7) + 2\,a(n-8) + a(n-10) - a(n-11)
\]
for every $n>0$.
\end{theorem}

\begin{proof}
The terms $a(0),\dots$ were computed exactly in integer arithmetic from the model, extended by
$A$ where the model itself was not evaluated, and $r(n)$ formed for each index. Those integers
vanish from $n=0+1$ onwards, and the run exceeds $S+11$, which by
Lemma~\ref{lem:crit} settles every larger index at once.
\end{proof}

\section{Verification}

Three checks, each independent of the algebra above.

First, the model reproduces all $53$ terms the entry publishes, exactly, in integer
arithmetic. This is what ties the model to the entry: it is built by reading the entry's
English, and a misreading gives different counts.

Second, the conjectured recurrence was evaluated directly on those published terms, with no
model involved, and holds wherever the proved range and the data overlap.

Third, the derived annihilator $A$ was itself tested on terms it had not been given: the model
was evaluated beyond the $S$ values that determine it and $A$ reproduced them. A bound that is
too small fails this test, which is why it is run.

\begin{thebibliography}{9}
\bibitem{oeis} The OEIS Foundation, \emph{The On-Line Encyclopedia of Integer
Sequences}, \url{https://oeis.org}, 2026. Sequence A183928.
\bibitem{beckrobins} M.~Beck and S.~Robins, \emph{Computing the Continuous Discretely},
2nd ed., Springer, 2015. (Ehrhart quasi-polynomials and their periods.)
\bibitem{stanley} R.~P.~Stanley, \emph{Enumerative Combinatorics}, Volume 1, 2nd ed.,
Cambridge University Press, 2011.
\bibitem{ds} F.~Diamond and J.~Shurman, \emph{A First Course in Modular Forms},
Springer GTM 228, 2005.
\end{thebibliography}

\end{document}
